% RESEARCH_DRAFT_SHA256: 01458905debf47e748a4fee1b9039d52c43306979ae32963cb399bdc0a9673de
% Layout and source-reference wrappers expanded for the web renderer.
\documentclass[11pt]{article}
\usepackage[margin=25mm]{geometry}
\usepackage{amsthm}
\usepackage{fontspec,unicode-math}
\setmainfont{Latin Modern Roman}
\setmathfont{Latin Modern Math}
\usepackage{amsmath,amssymb,mathtools,xurl,hyperref}
\hypersetup{colorlinks=true,urlcolor=blue,linkcolor=blue}
\setcounter{secnumdepth}{0}
\setlength{\parindent}{0pt}
\setlength{\parskip}{5pt}
\setlength{\emergencystretch}{3em}
\newcommand{\sref}[2]{\hyperlink{source:#1:#2}{[S#2]}}
\newcommand{\cataloguescope}[1]{\paragraph{Source record.}#1}
\newcommand{\floor}[1]{\left\lfloor#1\right\rfloor}
\newtheorem{proposition}{Proposition}
\newtheorem{lemma}[proposition]{Lemma}
\begin{document}
\paragraph{Research manuscript.} OpenAI \texttt{gpt-6-astra}, reasoning \texttt{ultra}; 11 October 2026.
\paragraph{Preserved source body.} The complete catalogue mathematical text and sources follow with layout and source links expanded, including source qualifications and its historical preparation notes. New research begins at the section headed Research attempt.

\section{Defining High-Dimensional Expanders Through Local Spectra}\label{statement}

\subsection{Definitions and mathematical statement}

\paragraph{Shared notation.}Standard number systems are $\mathbb N,\mathbb Z,\mathbb Q,\mathbb R,\mathbb C$. Specialized notation and index ranges follow the cited source and the definitions below.


In the LOCAL model vertices communicate synchronously along graph edges; messages and local computation are unrestricted and complexity is the number of rounds. After $t$ rounds a vertex can depend only on its radius-$t$ neighborhood and any supplied identifiers or randomness.

A finite simplicial complex is a finite family of vertex subsets closed under taking subsets. Its dimension is one less than its largest simplex size. A piecewise-linear embedding is an injective map linear on simplices after finite subdivisions.
\paragraph{Record scope and qualifications.}
The catalogue title is retained as metadata. Named certificates, contracts, bounds or frameworks in that title are not asserted as results of the original source extract.
The supplied registry flags this extraction as unrecoverable. It is displayed as a literal source fragment; the definitions provide context rather than a replacement formal statement.

Source-record limitation: Source flags the statement as unrecoverable: The official report uses this sentence as an expository heading, immediately supplies a local-spectral definition, and lists different concrete research directions afterward; it is not a distinct unresolved claim. No separate clean statement is promoted; the original record remains publicly visible for provenance.
The source text below is rendered literally where necessary; damaged formulas and composite extracts are not silently repaired into a new mathematical claim.
The supplied catalogue marks this record as unrecoverable. Its extract and terminology are documented, but a complete original problem has not been reconstructed.
\paragraph{Statement recorded in the supplied source.}
What is a high-dimensional expander; in particular, which local or spectral expansion conditions on the links of a bounded-dimensional simplicial complex should define this notion? \href{https://doi.org/10.4171/owr/2018/51}{[S2]}
\paragraph{Source status.}
The supplied catalogue labels this source record open; that metadata was checked against the supplied snapshot on 2026-10-11. The label is preserved for provenance and is not an independent verification of present open status or mathematical correctness. The catalogue literature assessment is dated 2026-08-21.
\subsection{Short English statement}
The supplied fragment concerns defining High-Dimensional Expanders Through Local Spectra. Its complete mathematical target is not recovered.
\subsection{Sources}
\begin{itemize}
\item[S1] Ulam AI and the UnsolvedMath Contributors, UnsolvedMath: A Curated Collection of Open Mathematics Problems, catalogue entry. Catalogue curation and metadata: CC BY 4.0. The source extract is retained; definitions, English summaries and qualifications are editorial additions.
\url{https://huggingface.co/datasets/ulamai/UnsolvedMath}.
\item[S2] Complexity Theory (2019), Oberwolfach Reports, DOI 10.4171/owr/2018/51, pp. 3025--3080
\url{https://doi.org/10.4171/owr/2018/51}.
\item[S3] Oberwolfach Report 51/2018
\url{https://ems.press/content/serial-article-files/46773}.
\item Source-specific Oberwolfach Reports licences, including CC BY-SA 4.0 where applicable, are retained; see THIRD\_PARTY\_NOTICES.md.
\end{itemize}
\paragraph{Preparation and review.}
The original source and UnsolvedMath contributors retain attribution for the source extract. Codex (GPT-6; exact serving revision not exposed) assisted with definitions, English summaries, scope notes and formatting. No new human mathematical authorship or independently verified solution is asserted. Catalogue identity checked against the complete supplied registry and current repository. Distinct upstream records, including repeated questions and broad or damaged extracts, are retained for complete coverage of the source CS category. A repeated record is not a new independent conjecture. Complete mathematical-equivalence and current-literature review has not been performed. Preparation and final review follow contribution prompts C and D, with source gaps disclosed. Structural and rendering checks do not establish a new theorem or open status.
\label{end:statement}

\section{Research attempt}
\subsection{Literature, source context and scope}
The full original catalogue body and source qualifications are preserved above. The relevant primary workshop passages were read in the cached publisher report. Historical claims and subsequent attributed results are distinguished from the mathematical work below. A restricted example, standard lemma or known theorem is not counted as new progress on the entire original target. Missing assumptions and repeated or rhetorical extracts are not silently turned into different solved problems.

\subsection{Why the source definition retains global connectedness}
The preserved record is flagged as an unrecoverable problem extraction. The primary page already gives a spectral definition and separately lists research directions; we do not invent a new unresolved definition problem. We test what goes wrong if one drops its explicit connectedness requirement and checks only nonempty-face links.
Let $r\ge4$ and take the disjoint union of two complete two-dimensional complexes on $r$ vertices each: each contains every vertex, edge and triangle on its vertex set, with no higher faces.
\begin{proposition}
Every vertex link has normalized adjacency spectrum $1,-1/(r-2),\ldots,-1/(r-2)$, whereas the underlying graph has eigenvalue 1 with multiplicity two and is disconnected.
\end{proposition}
\begin{proof}
The link of a vertex consists of every edge between its other $r-1$ component vertices, so its graph is $K_{r-1}$. The adjacency matrix of $K_t$ is $J-I$. After dividing by degree $t-1$, it acts as 1 on the all-ones vector and as $-1/(t-1)$ on the orthogonal subspace where $J$ vanishes. Substituting $t=r-1$ gives the link spectrum. The underlying graph is the disjoint union of two copies of $K_r$, whose normalized adjacency matrix is block diagonal. Each component supplies an independent eigenvector of eigenvalue 1, and there are no edges between components.
\end{proof}
Thus all vertex-link nontrivial eigenvalues can have arbitrarily small absolute value while the global graph is disconnected. In dimension two these are precisely the nonempty faces whose links have a graph of positive dimension; edge links are zero-dimensional. The example fails the source's first connectedness condition, so it is not a counterexample to the definition or to a properly hypothesized local-to-global theorem.
\subsection{Verification and remaining gaps}
The complexes and link graphs were enumerated for $r=4,5,6$, and their spectra checked. This supplies a concrete reason that a tempting weakened definition is inadequate. It does not choose among all notions of high-dimensional expansion, settle the source's later coding/PCP questions, or turn an expository heading into a new theorem target. No original-target completion is claimed.

\section{References and source attribution}
Primary workshop report: \url{https://doi.org/10.4171/owr/2018/51}. Relevant passages read: 2018/51 pp. 3041--3043, expository heading, explicit connectedness/local-spectrum definition and separate research directions.

All retrievals were checked on 11 October 2026. Original question attribution remains in the preserved Sources section. External theorems are attributed and not claimed as new discoveries or fully verified merely by citation.


\paragraph{Integration review.} Codex (GPT-6; exact serving revision not exposed) reviewed the written arguments and expanded layout and source-reference wrappers. Repository metadata and the production web renderer were checked using Python 3.13 and JavaScript with MathJax 3.2.2. This model-assisted review is not independent human proof verification; no PDF compilation or formal proof checking is asserted.

\section{Completion Estimate}
\noindent\textbf{COMPLETION: 0\%}
No substantiated new progress toward the entire original target is claimed. The worked results and failed route are delimited above; a known result, illustrative model, or auxiliary obstruction is not counted as original-target completion.

\section{Final status and attribution}
\textbf{Attempt status: unresolved.} This describes this research attempt, not a new ruling on the source catalogue's status or on external literature claims.
The mathematics and this self-audit were generated by OpenAI \texttt{gpt-6-astra}, reasoning \texttt{ultra}, on 11 October 2026. Codex (GPT-6; exact serving revision not exposed) orchestrated the batch, reviewed the written mathematical arguments, and contributed formatting/repository review. This additional model review is not independent human verification. No human mathematical authorship, independent human proof checking, or formal verification is claimed. The historical preparation paragraphs in the reproduced catalogue body concern that body only.

Contribution prompt A was applied to the complete definition and supplied source record, followed by the adversarial checks of prompt B. Local reproduction files are \texttt{batch022/checks.py}, \texttt{computations.json}, and the per-ID audit JSON. They were not public links at preparation time. Numerical checks are finite and do not prove unrestricted results. No LaTeX engine was found locally; compilation was not performed by this research agent.

\end{document}
